\documentclass[11pt]{article}
\usepackage[T1]{fontenc}
\usepackage{lmodern}
\usepackage{amsmath,amssymb}
\usepackage{booktabs}
\usepackage{array}
\usepackage[a4paper,margin=2.5cm]{geometry}
\usepackage[hidelinks]{hyperref}
\newcommand{\doi}[1]{\href{https://doi.org/#1}{\texttt{#1}}}
\setlength{\parskip}{4pt}
\title{\vspace{-1.6cm}\Large Guided Lamb Waves in a Piezoelectric Semiconductor Plate with a Size Effect and a\\[2pt]
Relaxation-Time Heat Equation\\[8pt]\large Research plan, first half}
\author{Dr.\ Vipin \\ Gurugram University}
\date{30 October 2026}
\begin{document}
\maketitle
\vspace{-6mm}

\section{Title, abstract, keywords}

\paragraph{Working title} Guided Lamb waves in a piezoelectric semiconductor plate: attenuation from carrier
transport, nonlocal strain gradient elasticity and a single-lag heat equation.

\paragraph{Abstract} We study plane-strain Lamb waves in a plate of thickness $2h$ cut from a piezoelectric
semiconductor of the ZnO family, with free faces at $z=\pm h$ and poling along the thickness. The model
carries five fields: the two displacements, the electric potential, the excess carrier density and the
temperature rise. The mechanical stress is written in the nonlocal strain gradient form, with one length
for the nonlocal softening and one for the gradient stiffening, so that the size effect is described by
two numbers that enter only as length over thickness. The heat equation is of the Moore-Gibson-Thompson (MGT)
type in its base form, with one phase lag and one time derivative cancelled, so that the
parabolic and the hyperbolic descriptions are recovered as limits. The carrier equation carries
diffusion, drift, a thermodiffusion term and an optional recombination term, and the heat equation carries
the reciprocal Peltier term, the thermoelastic coupling and the pyroelectric coupling. Three questions are
asked. How much attenuation does the carrier field add to each Lamb branch. Over what thickness does a
single relaxation time change the attenuation by more than the experimental resolution. How large must the
two size lengths be for the size effect to be visible in the dispersion at all. The plan describes the
model, the sequence of comparisons that will be made before any new result is claimed, and the four
figures the paper will carry. Numbers are deliberately absent at this stage.

\paragraph{Keywords} guided waves; piezoelectric semiconductor; nonlocal strain gradient elasticity;
thermoelastic coupling; carrier transport; attenuation.

\section{Introduction}

Piezoelectric semiconductor plates are the working element of a large family of acoustic and optical
devices, and the guided waves in them are used both to sense and to filter. The mechanics of such a plate
differs from that of a dielectric plate in one respect that matters here: the electric field is not free to
arrange itself as it likes, because mobile carriers redistribute under the field it produces. Yang and Zhou
solved the amplification of acoustic waves in piezoelectric semiconductor plates with the drift-diffusion
pair that most later work follows \cite{YangZhou2005}, and guided waves in semiconducting plates of graded
or layered cross section were taken up with numerical machinery that can carry a variable coefficient
through the thickness \cite{Othmani2017}. In all of that work the plate is isothermal. The temperature is
either absent or treated as a constant that shifts the material constants.

The thermal side of our problem is already partly occupied, and the plan says so at the start. Gupta and
co-workers put the MGT heat equation, with one lag, into a piezo-semiconductor half space with
voids and studied the size-dependent Rayleigh wave \cite{Gupta2025}. That paper settles two things for us:
the heat operator in a semiconducting piezoelectric medium is writable, and the Rayleigh limit of a
half space is already published. What it does not contain is a finite plate, the free-face problem, the
nonlocal strain gradient description of the elastic stress, or the reciprocity between the thermodiffusion
term in the carrier equation and the Peltier term in the heat equation. Our contribution has to be read as
that list and nothing wider.

A second body of work puts the temperature back, but into a dielectric. Sharma, Pal and Chand solved the
piezothermoelastic plate with free, electrically shorted and thermally insulated faces, obtained closed
secular equations for symmetric and antisymmetric families, and reported attenuation coefficients for
cadmium selenide \cite{Sharma2004}. Kumar and Kansal added a diffusing species to a transversely isotropic
thermoelastic plate and obtained the same kind of closed form \cite{KumarKansal2008}. In both, the third
field is a concentration that does not carry charge, so the electric and diffusive problems do not talk to
each other in the way they do in a semiconductor.

The two lines have recently met in the continuum theory of Qu, Pan, Zhu, Jin and Roy, who derive a
thermoelectric theory of piezoelectric semiconductors with the Seebeck and Peltier effects built in from a
dissipation inequality \cite{Qu2023}. That paper settles the constitutive side of our problem, and we
should not pretend otherwise. What it does not do is the plate: no faces, no guided modes, no dispersion
relation, no attenuation curve against frequency times thickness.

The size effect is a separate axis. Lim, Zhang and Reddy combined the nonlocal integral law with strain
gradient elasticity so that one length softens and another stiffens, and gave the equivalent differential
form together with the constitutive boundary conditions that come with it \cite{Lim2015}. Barretta and
Marotti de Sciarra examined what those boundary conditions do in bounded domains and showed that the
differential law raises the order of the problem \cite{Barretta2018}. For a plate this is not a detail:
the order in the thickness coordinate goes up, and the faces must be given more conditions than the
classical six.

For the thermal side we use the MGT theory in its base form, one phase lag, one time
derivative cancelled, which is the version implemented in the recent heat transfer literature
\cite{Gupta2025}. The operator is simple enough to carry through a dispersion calculation and it has two
limits that can be tested against the classical solutions.

The gap, then, is a finite plate in which all four ingredients sit together: a semiconducting carrier
field, a temperature field with a relaxation time, a mechanical size effect, and free faces. To our
knowledge, after searching Crossref and the open web in October 2026 for guided waves in piezoelectric
semiconductor plates, nobody has solved that. The honest version of the claim is that our contribution is
the free-face problem, the size effect and the attenuation curves, not any new constitutive law. The
thermoelectric constitutive relations are Qu et al.'s, the heat operator in a semiconductor is Gupta et
al.'s, the nonlocal strain gradient form is Lim et al.'s, and we should say so in the abstract of the final
paper.

Three things follow from that reading, and they shape the plan. The carrier coupling is small, because
both the thermodiffusion coefficient and the Peltier coefficient are proportional to the carrier density,
and the carrier density in a good piezoelectric is not large. The thermal relaxation is visible only in a
band of thickness, because the relaxation time multiplies the square of the thermal diffusion length. The
size effect enters only as the ratio of a length to the thickness, so it does nothing until the plate is
thin. A paper that puts all three effects in and then reports that they are all negligible would be
correct and useless. The plan therefore fixes, before any computation, the test that decides whether each
effect is worth a figure.

\section{Problem statement and assumptions}

We take a plate of thickness $2h$, infinite in $x$ and $y$, with free faces at $z=\pm h$. Motion is plane
strain in the $x$--$z$ plane, and the poling direction is $z$, so the material is transversely isotropic
with the axis along the thickness. Fields vary as $\exp[i(kx-\omega t)]$ with $k=k_r+ik_i$, and a positive
$k_i$ means decay along $+x$. The frequency is reported through the product $f\,(2h)$, so that plates of
different thickness can be compared on one axis. The five unknowns are $u_x$, $u_z$, $\varphi$, $\Delta n$
and $\theta$, with $\Delta n$ the excess carrier density over equilibrium and $\theta=T-T_0$ the
temperature rise.

The assumptions are these, and each one is a place where a referee can push.

\begin{enumerate}
\item Linear piezoelectricity at constant temperature, with the elastic, piezoelectric and dielectric
constants taken at $T_0$ and not varied with $\theta$.
\item Small carrier perturbation, $\Delta n\ll n_0$, so that the drift term can be linearised about the
bias field and the recombination term is linear in $\Delta n$.
\item A uniform bias field $E_0$ along $x$ that produces the drift velocity $v_d=b_nE_0$; the bias field is
not solved for, it is an input.
\item No photoexcitation, no trapping, no surface charge on the faces beyond what the electrical condition
imposes.
\item Small temperature rise, $|\theta|/T_0\ll1$, so the heat equation can be written for $\theta$ and the
coefficients can be held constant.
\item Homogeneous, stress-free bulk material with no initial stress and no piezoelectric stiffening from a
pre-load.
\end{enumerate}

Two consequences of the second assumption and the fifth should be stated at the outset. The carrier term in
the heat equation is reversible in the thermodynamic sense and is not a dissipation; the dissipation comes
from the Joule and diffusion terms, which we keep in the entropy balance but do not put in the field
equations above linear order. And the size lengths $l_{\rm nl}$ and $l_g$ are material inputs, not fitting
parameters. If the dispersion is insensitive to them at the thicknesses of interest, that is a result and
it goes in the paper.

\section{The model}

\subsection{Constitutive relations}

Strain and field are $\varepsilon_{ij}=\tfrac12(u_{i,j}+u_{j,i})$, $E_i=-\varphi_{,i}$, with $i,j=x,z$.
The total stress carries the nonlocal operator, the local stress carries the gradient operator on the
elastic part only,
\begin{equation}
\bigl(1-l_{\rm nl}^2\nabla^2\bigr)t_{ij}=\sigma_{ij},\qquad
\sigma_{ij}=c_{ijkl}\bigl(1-l_g^2\nabla^2\bigr)\varepsilon_{kl}-e_{kij}E_k-\beta_{ij}\theta ,
\end{equation}
where the differential form of the nonlocal strain gradient law is that of Lim, Zhang and Reddy \cite{Lim2015} (their Eqs.\ (17)--(19), as we have them quoted in \cite[Eq.\ (18)]{Barretta2018}; open Section~2 of \cite{Lim2015} and compare, we have not seen the page ourselves). The elastic part with the operator $(1-l_g^2\nabla^2)$ is the strain-gradient term of Mindlin \cite{Mindlin1964}; the piezoelectric part of $\sigma_{ij}$ follows \cite[Eq.\ (1)]{Zhu2023} and the thermal stress follows \cite[Eq.\ (3.1)]{Quintanilla2019}. The electric displacement is local,
\begin{equation}
D_i=e_{ikl}\varepsilon_{kl}+\epsilon_{ik}E_k+p_i\theta ,
\end{equation}
with $\beta_{ij}=c_{ijkl}\alpha_{kl}$ from the thermal expansion tensor and $p_i$ the pyroelectric vector. The first two terms are \cite[Eq.\ (1)]{Zhu2023}. The strain--displacement and field--potential relations above are \cite[Eq.\ (2)]{Zhu2023}. The term $p_i\theta$ is the standard pyroelectric addition of thermopiezoelectricity \cite{Mindlin1974}; we did not read its equation number, so please locate it there.
The choice to keep the piezoelectric and pyroelectric terms local is deliberate. Only the elastic strain
energy is allowed to feel the two lengths, so the electric problem stays second order, no extra electrical
condition is needed at the faces, and no term that is already a field gradient is differentiated twice.

\subsection{Field equations}

Momentum. The balance law is $\rho\ddot u_i=t_{ij,j}$ \cite[Eq.\ (3.3)]{Quintanilla2019}, in the static form $\sigma_{ij,j}=0$ also \cite[Eq.\ (4)]{Zhu2023}. Applying the operator of Eq.~(1) to it gives the equivalent form in which the nonlocal operator acts on the inertia; this step is ours and you should redo it,
\begin{equation}
\bigl(1-l_{\rm nl}^2\nabla^2\bigr)\rho\,\ddot u_i=\sigma_{ij,j},
\end{equation}
Gauss law with the mobile charge, \cite[Eqs.\ (4)--(6)]{Zhu2023} (there with both carrier types and the doping terms; we keep one carrier type, electrons of charge $-q$, and the doping cancels against $n_0$),
\begin{equation}
D_{i,i}=-q\,\Delta n ,
\end{equation}
carrier continuity. The drift--diffusion current of the electrons is \cite[Eq.\ (3)]{Zhu2023}, linearised as in \cite[Eq.\ (7)]{Zhu2023}, and its divergence enters the continuity law; the time-dependent form is that of \cite{YangZhou2005} (equation numbers not read by us). The drift term with $v_d$, the thermodiffusion term and the recombination term are our additions, each following from a textbook law,
\begin{equation}
\frac{\partial\Delta n}{\partial t}
=D_n\nabla^2\Delta n+b_nn_0\nabla\!\cdot\!E+v_d\frac{\partial\Delta n}{\partial x}
+\beta_{\rm th}\nabla^2\theta-\frac{\Delta n}{\tau_R},
\end{equation}
and the heat equation in the base form of the Moore--Gibson--Thompson theory. The MGT heat equation is \cite[Eq.\ (2.5)]{Quintanilla2019} and its thermoelastic version is \cite[Eq.\ (3.5)]{Quintanilla2019}; with the rate conductivity $k^*$ put to zero and one time integration (zero initial data) this becomes the equation below (the source sets $T_0=1$, so restore $T_0$ in the coupling term). The terms with $p_i$, $\Pi_{\rm th}$ and $E_g$ are ours,
\begin{equation}
\Bigl(1+\tau_q\frac{\partial}{\partial t}\Bigr)
\Bigl[\rho c_p\dot\theta+T_0\beta_{ij}\dot\varepsilon_{ij}+T_0p_i\dot E_i
+q\Pi_{\rm th}\frac{\partial\Delta n}{\partial t}+\frac{E_g\Delta n}{\tau_R}\Bigr]
=k_{\rm th}\nabla^2\theta .
\end{equation}
Five equations for five unknowns. The two coefficients that appear in both the carrier and the heat
equation are not independent. With $S_n$ the Seebeck coefficient of the electrons,
\begin{equation}
\beta_{\rm th}=-b_nn_0S_n,\qquad \Pi_{\rm th}=T_0S_n,\qquad
\frac{\Pi_{\rm th}}{\beta_{\rm th}}=-\frac{T_0}{b_nn_0},
\end{equation}
which is the Onsager symmetry between the Soret and Peltier cross effects (a textbook relation of irreversible thermodynamics; the semiconductor continuum of \cite{Qu2023} is built on the same relations, equation numbers not read), and the Kelvin relation
$\Pi=TS$ in its local form. For $n$-type material the carriers have charge $-q$ and $S_n<0$, so
$\beta_{\rm th}>0$ and $\Pi_{\rm th}<0$. Check these two signs yourself against the physical picture, that is, which way the electrons move under a temperature gradient and where the heat is released, before you use them. The recombination pair $-\Delta n/\tau_R$ and
$+E_g\Delta n/\tau_R$ is an option, not a commitment; taking $\tau_R\to\infty$ deletes both terms and
nothing else changes. If you keep it, you must justify a value for $\tau_R$ from the literature or sweep it
and say so.

\subsection{Notation}

\begin{center}
\small
\begin{tabular}{p{2.4cm}p{6.7cm}p{3.4cm}}
\toprule
symbol & meaning & unit\\
\midrule
$u_x,u_z$ & displacements & m\\
$\varphi$ & electric potential & V\\
$\Delta n$ & excess carrier density & m$^{-3}$\\
$\theta$ & temperature rise over $T_0$ & K\\
$k=k_r+ik_i$ & wavenumber; $k_i>0$ is decay along $+x$ & m$^{-1}$\\
$2h$ & plate thickness; faces at $z=\pm h$ & m\\
$c_{ijkl}$, $e_{kij}$, $\epsilon_{ik}$ & elastic, piezoelectric, dielectric constants & Pa, C\,m$^{-2}$, F\,m$^{-1}$\\
$\rho$ & mass density & kg\,m$^{-3}$\\
$l_{\rm nl}$, $l_g$ & nonlocal and gradient lengths, used as $l/h$ & m\\
$q$ & magnitude of the carrier charge; electron charge is $-q$ & C\\
$n_0$ & equilibrium carrier density & m$^{-3}$\\
$b_n$ & electron mobility & m$^2$V$^{-1}$s$^{-1}$\\
$D_n$ & carrier diffusion coefficient & m$^2$s$^{-1}$\\
$E_0$, $v_d$ & bias field, drift velocity $v_d=b_nE_0$ & V\,m$^{-1}$, m\,s$^{-1}$\\
$S_n$ & Seebeck coefficient of the electrons, negative for $n$-type & V\,K$^{-1}$\\
$\beta_{\rm th}$ & thermodiffusion coefficient, $-b_nn_0S_n$, positive for $n$-type & m$^{-1}$s$^{-1}$K$^{-1}$\\
$\Pi_{\rm th}$ & Peltier coefficient, $T_0S_n$ & V\\
$\tau_R$, $E_g$ & recombination time, band gap & s, J\\
$\alpha_{kl}$, $\beta_{ij}$, $p_i$ & thermal expansion, thermal stress $c_{ijkl}\alpha_{kl}$, pyroelectric vector & K$^{-1}$, Pa\,K$^{-1}$, C\,m$^{-2}$K$^{-1}$\\
$\rho c_p$, $k_{\rm th}$, $\kappa$ & heat capacity per volume, conductivity, diffusivity $k_{\rm th}/(\rho c_p)$ & J\,m$^{-3}$K$^{-1}$, W\,m$^{-1}$K$^{-1}$, m$^2$s$^{-1}$\\
$\tau_q$, $\tau_T$ & phase lag, $\tau_T=(2h)^2/(\pi^2\kappa)$ & s\\
$\omega$, $f$ & angular frequency, frequency & rad\,s$^{-1}$, Hz\\
\bottomrule
\end{tabular}
\end{center}

\subsection{Order of the system, and where this chapter stops}

The nonlocal and gradient operators raise the order of the problem in $z$, so with free faces the plate
needs more than the classical traction conditions, and the extra conditions have to be chosen with care
because the differential law is only equivalent to the integral one when they are consistent with it. Which
extra conditions to apply at $z=\pm h$, and how to close the electrical and thermal conditions there, is
your problem to settle; this chapter stops before the faces.

\section{Work plan}

Six steps, in this order. Each step ends with a written comparison and a number, and no step is allowed to
borrow a number from the next.

\paragraph{Step 1. Elastic plate, no fields switched on} With $e=0$, $\Delta n=0$, $\theta=0$, $l_{\rm nl}=0$
and $l_g=0$, the model must return the classical Lamb spectrum of an isotropic plate, both families, and
the Rayleigh velocity in the short wavelength limit. Compare against the closed-form Rayleigh value
computed from the same $\lambda$ and $\mu$ and require agreement to a relative $10^{-3}$ at large
frequency times thickness. This is the floor of the whole paper; if it fails nothing else is worth
computing.

\paragraph{Step 2. Piezoelectric stiffening alone} Switch on $e_{kij}$ and keep the electrical faces
consistent with the sheet condition you choose. The dispersion must split the branches in the way the
published piezoelectric plate results do, and the modal coupling must vanish for the family that
decouples. Compare the shape of the branches, not the numbers, since our material set is not the one in
those papers.

\paragraph{Step 3. Heat alone, in the two limits} Switch on $\theta$ with $\Delta n=0$ and compare the two
limits of the heat equation: the isothermal limit, in which the temperature follows the strain instantly,
and the adiabatic limit, in which no heat flows. The two must bracket the coupled answer to a width of
order the thermoelastic coupling parameter. Report the width. This bracket is the test of the model, not
the result of the paper.

\paragraph{Step 4. Thermal relaxation switched on} With $\tau_q$ finite the attenuation must saturate at
large frequency and the peak must sit at $1/|\tau_T-\tau_q|$ rather than at $1/\tau_T$, with the plateau
rescaled by $\tau_T^2/(\tau_T-\tau_q)^2$. Confirm both scalings against a one-dimensional reduction of the
same equations before the plate is solved. Then draw the visibility map: thickness on one axis, $\tau_q$ on
the other, and shade the region where the shift of the peak exceeds the linewidth of the branch.

\paragraph{Step 5. Carriers switched on, and the switching-off test} Set $\beta_{\rm th}=0$ and $\Pi_{\rm th}=0$.
The carrier field must still change the dispersion through Gauss law alone, since the charge enters
$D_{i,i}$, and this is a test of the assembly rather than of the physics. Then the reverse: keep the
carrier terms and set the mechanical coupling $e_{kij}=0$. The result must be a purely electrical
relaxation with no acoustic branch shift. In both cases the comparison is numerical: the difference between
the full model with a coupling set to zero and the reduced model without that coupling must be below
$10^{-6}$ in the relative sense on $k_r$ and $k_i$ for every sampled mode. Log the residuals in a table in
the appendix. A test that passes at $10^{-4}$ is a warning, not a pass.

\paragraph{Step 6. Size effect and the final sweep} Vary $l_{\rm nl}/h$ and $l_g/h$ over a range that spans
the values quoted for the material, keeping their sum fixed so that softening and stiffening are separated
from the net shift. Then produce the two attenuation families against frequency times thickness, for three thicknesses. When the thermodiffusion and Peltier terms are on, the symmetric and antisymmetric families no longer separate exactly (the reason is in the third file of this package), so label each branch by the displacement pattern it follows to the elastic limit. Report the carrier contribution as a difference of
attenuations, so that a reader who does not care about the carriers can still use the plate results.

\section{Figures, risks and timeline}

The paper will carry four figures and no more. First, the dispersion of the elastic limit and of the full
model on the same axes, so that the shift is visible without a second panel. Second, the attenuation of the
two lowest branches against frequency times thickness, with the isothermal and adiabatic brackets drawn as
bands. Third, the visibility map of Step 4, with the usable thickness range shaded. Fourth, the
carrier contribution to attenuation as a function of drift velocity, for two carrier densities, since the
drift velocity is the one knob a device engineer can actually turn.

The risks, in order of how much damage they do. The first is identifiability: the thermal and carrier
contributions to attenuation may be of the same size as the spread between the two size lengths, in which
case no experiment can separate them and the paper must say so and present a sensitivity analysis instead
of a fit. The second is that the relaxation time is not resolvable at the thicknesses where the continuum
heat model is trustworthy; the map of Step 4 is designed to expose that rather than hide it. The third is
the boundary conditions of the size effect, which are the subject of an active argument in the literature
\cite{Barretta2018} and which we will not settle. The fourth is the material set: the piezoelectric
constants of ZnO are quoted over a range that is wider than the effect we are chasing, and we must show
that every conclusion survives the range.

Timeline, twelve weeks from the start of work. Weeks one to two, Step 1 and Step 2, and reading of the four
papers that decide the claim \cite{Qu2023,Sharma2004,Othmani2017,YangZhou2005}. Weeks three to four, Step 3
and the thermal side, including the one-dimensional reduction of Step 4. Weeks five to six, Step 4 and the
visibility map. Weeks seven to eight, Step 5, and the switching-off log. Weeks nine to ten, Step 6 and the
sweeps. Weeks eleven to twelve, figures, writing, and the internal read-through before the draft goes to
me.

\begin{thebibliography}{99}
\small
\bibitem{YangZhou2005} J.~S.~Yang, H.~G.~Zhou, \emph{Int.\ J.\ Solids Struct.} \textbf{42} (2005) 3171, \doi{10.1016/j.ijsolstr.2004.10.011}.
\bibitem{Othmani2017} C.~Othmani, F.~Takali, A.~Njeh, \emph{Superlattices Microstruct.} \textbf{106} (2017) 86, \doi{10.1016/j.spmi.2017.03.036}.
\bibitem{Sharma2004} J.~N.~Sharma, M.~Pal, D.~Chand, \emph{J.\ Therm.\ Stresses} \textbf{27} (2004) 33, \doi{10.1080/01495730490255709}.
\bibitem{KumarKansal2008} R.~Kumar, T.~Kansal, \emph{Int.\ J.\ Solids Struct.} \textbf{45} (2008) 5890, \doi{10.1016/j.ijsolstr.2008.07.005}.
\bibitem{Qu2023} Y.~Qu, E.~Pan, F.~Zhu, F.~Jin, A.~K.~Roy, \emph{Int.\ J.\ Eng.\ Sci.} \textbf{182} (2023) 103775, \doi{10.1016/j.ijengsci.2022.103775}.
\bibitem{Lim2015} C.~W.~Lim, G.~Zhang, J.~N.~Reddy, \emph{J.\ Mech.\ Phys.\ Solids} \textbf{78} (2015) 298, \doi{10.1016/j.jmps.2015.02.001}.
\bibitem{Barretta2018} R.~Barretta, F.~Marotti de Sciarra, \emph{Int.\ J.\ Eng.\ Sci.} \textbf{130} (2018) 187, \doi{10.1016/j.ijengsci.2018.05.009}.
\bibitem{Gupta2025} V.~Gupta, F.~A.~Awwad, E.~A.~A.~Ismail, H.~Ahmad, S.~Das, \emph{Int.\ J.\ Numer.\ Methods Heat Fluid Flow} \textbf{35}(9) (2025) 3400, \doi{10.1108/HFF-04-2025-0254}.
\bibitem{Krahl2020} F.~Krahl, A.~Giri, M.~S.~B.~Hoque, L.~Sederholm, P.~E.~Hopkins, M.~Karppinen, \emph{J.\ Phys.\ Chem.\ C} \textbf{124} (2020) 24731, \doi{10.1021/acs.jpcc.0c06461}.
\bibitem{Ozgur2005} \"U.~\"Ozg\"ur et al., \emph{J.\ Appl.\ Phys.} \textbf{98} (2005) 041301, \doi{10.1063/1.1992666}.
\bibitem{Zhu2023} J.~Zhu, M.~Negahban, J.~Xu, R.~Xia, Z.~Li, \emph{Micromachines} \textbf{14} (2023) 2174, \doi{10.3390/mi14122174}.
\bibitem{Quintanilla2019} R.~Quintanilla, \emph{Math.\ Mech.\ Solids} \textbf{24} (2019) 4020, \doi{10.1177/1081286519862007}.
\bibitem{Charal2019} A.~Charalambopoulos, E.~Douka, M.~Mavratzas, Second order formulation of boundary value problems in gradient elasticity, arXiv:1909.10776 (2019).
\bibitem{Mindlin1964} R.~D.~Mindlin, \emph{Arch.\ Rational Mech.\ Anal.} \textbf{16} (1964) 51.
\bibitem{Mindlin1974} R.~D.~Mindlin, \emph{Int.\ J.\ Solids Struct.} \textbf{10} (1974) 625.
\end{thebibliography}\end{document}
